% =========================================================================
% Paper Clay Closure Perturbative : Two-input perturbative closure of the
%        SU(N) Yang--Mills mass gap chain on T^4 (regime epsilon <= c_0/sqrt(beta))
%
% Author : Kevin Remondiere (Independent Researcher, Oloron-Sainte-Marie)
% Date   : 26 May 2026
% License: CC-BY-4.0
% Target : Communications in Mathematical Physics
% =========================================================================

\documentclass[11pt,a4paper]{amsart}

\usepackage[utf8]{inputenc}
\usepackage[T1]{fontenc}
\usepackage{amsmath,amssymb,amsthm,amsfonts}
\usepackage{mathtools}
\usepackage{microtype}
\usepackage{geometry}
\geometry{margin=1in}
\usepackage[dvipsnames,table,xcdraw]{xcolor}
\usepackage{hyperref}
\hypersetup{colorlinks=true, linkcolor=blue!50!black, citecolor=blue!50!black, urlcolor=blue!50!black}
\usepackage{booktabs}
\usepackage{enumitem}

\newtheorem{theorem}{Theorem}[section]
\newtheorem{lemma}[theorem]{Lemma}
\newtheorem{proposition}[theorem]{Proposition}
\newtheorem{corollary}[theorem]{Corollary}
\theoremstyle{definition}
\newtheorem{definition}[theorem]{Definition}
\newtheorem{remark}[theorem]{Remark}
\newtheorem*{notation*}{Notation}

\newcommand{\R}{\mathbb{R}}
\newcommand{\Z}{\mathbb{Z}}
\newcommand{\C}{\mathbb{C}}
\newcommand{\gfrak}{\mathfrak{g}}
\newcommand{\hfrak}{\mathfrak{h}}
\newcommand{\su}{\mathfrak{su}}
\newcommand{\norm}[1]{\left\lVert#1\right\rVert}
\newcommand{\abs}[1]{\left\lvert#1\right\rvert}
\newcommand{\inner}[2]{\langle #1, #2\rangle}
\newcommand{\Ric}{\mathrm{Ric}}
\newcommand{\Tr}{\mathrm{Tr}}
\newcommand{\ad}{\mathrm{ad}}
\newcommand{\Coul}{\mathrm{Coul}}
\newcommand{\Hess}{\mathrm{Hess}}
\newcommand{\dvol}{\, d\mathrm{vol}}
\newcommand{\kFP}{\kappa_{\mathrm{FP}}}
\DeclareMathOperator{\spec}{spec}
\DeclareMathOperator{\rk}{rk}
\DeclareMathOperator{\id}{Id}
\DeclareMathOperator{\sym}{Sym}

\title[Two-input perturbative closure on $T^4$]{Two-Input Perturbative Closure of the $SU(N)$ Yang--Mills Mass Gap Chain on $T^4$: Polchinski Cancellation and Zegarlinski--Gribov Decomposition}
\author{K\'evin R\'emondi\`ere}
\address{Independent Researcher, Oloron-Sainte-Marie, France}
\email{kevin.remondiere@gmail.com}
\thanks{ORCID: \href{https://orcid.org/0009-0008-2443-7166}{0009-0008-2443-7166}.}
\date{May 26, 2026}
\subjclass[2020]{Primary 81T13, 81T17; Secondary 58J35, 60H30, 60J60, 39B62.}
\keywords{Yang--Mills mass gap, Polchinski equation, Faddeev--Popov determinant, structure constants antisymmetry, $d$-symbol, Zegarlinski decomposition, Gribov horizon, logarithmic Sobolev inequality, lattice gauge theory.}

\begin{document}

\begin{abstract}
Building on the explicit Uniform FP Hessian Bound established in
\cite{RemondiereFPHess_2026} and on the conditional spectral framework
\cite{RemondiereKRFP3_2026, RemondiereKRFPBakryEmery_2026}, we attempt
to discharge the two standard analytic inputs that were identified as
the remaining obstructions for an unconditional perturbative-regime mass
gap on $T^4$: (i) the Polchinski cubic-term cancellation
$\langle \nabla V_t, C_t \cdot \nabla \Hess V_t\rangle = 0$ for the
$SU(N)$ Wilson effective potential, and (ii) the Zegarlinski-style
good-set/bad-set decomposition compatible with the Gribov horizon
$\partial \Omega$ of \cite{Gribov1978, DellAntonioZwanziger1991}.

We prove:
\begin{itemize}[leftmargin=2em]
\item \textbf{Theorem~\ref{thm:cancel-leading}} (Cancellation at the leading
non-trivial order). For the bare Wilson potential $V_0 = \beta S_W -
\log\det M[A]$ on $SU(N)$ in Coulomb gauge, the cubic term
$\langle \nabla V_0, C_0 \cdot \nabla \Hess V_0\rangle$ vanishes
\emph{exactly} at order $A^0$ (vacuum) and is bounded by
$O(N^4 g^4 \varepsilon^2)$ at order $A^2$, provided the $d$-symbol
contractions are estimated via the explicit Lie-algebraic identity of
Lemma~\ref{lem:dsym-contraction}.
\item \textbf{Theorem~\ref{thm:cancel-flow}} (Conditional all-order
cancellation along the flow). Assuming a structural hypothesis
(\textsc{Hyp-CST}) on the conservation of total $f$-antisymmetry under the
Polchinski flow, the cubic term vanishes for all $t \geq 0$ in the
perturbative regime $\norm{A}_{L^\infty} \leq \varepsilon \leq
c_0/(C_1 \sqrt{N\beta})$. We isolate (\textsc{Hyp-CST}) precisely and
honestly identify the residual gap.
\item \textbf{Theorem~\ref{thm:zeg-gribov}} (Zegarlinski--Gribov
decomposition, conditional). For the gauge-fixed $SU(N)$ Wilson measure
$\mu_\beta$ on $T^4_{L,a}$, the decomposition
$\Omega = \Omega_{\mathrm{good}} \cup \Omega_{\mathrm{bad}} \cup
\partial \Omega$ satisfies:
$\mu_\beta(\Omega_{\mathrm{bad}}) \leq C e^{-c \beta L^4}$ (standard
Driver-type concentration \cite{Driver1989}), $\mu_\beta(\partial \Omega) = 0$
(Gribov horizon $= \{\det M[A] = 0\}$, FP determinant in the numerator),
and on $\Omega_{\mathrm{good}}$ the LSI constant satisfies
$C_{\mathrm{LSI}}(\mu_\beta|_{\Omega_{\mathrm{good}}}) \leq C_0(N,D)$
uniform in $L$, conditional on
(\textsc{Hyp-CST}) and the standard Zegarlinski recovery lemma of
\cite{Zegarlinski1996, BakryGentilLedoux2014}.
\end{itemize}

\medskip
\noindent\textbf{Honest scope.} Theorem~\ref{thm:cancel-leading} is
proved unconditionally at leading order. Theorem~\ref{thm:cancel-flow}
\emph{requires} the explicit hypothesis (\textsc{Hyp-CST}), stating that
the $d$-symbol contributions to $\nabla \Hess V_t$ remain
sub-leading along the flow. We discuss this hypothesis honestly in
\S\ref{sec:Hyp-CST}, identify two routes to proving it
(one via a Brydges-Federbush-type cluster expansion, one via the
Bauerschmidt--Dagallier $\varphi^4_3$-style estimates adapted to $SU(N)$
\cite{BauerschmidtDagallier2024}), and provide a probability estimate
$50$--$70\%$ over $3$--$6$ months for a complete proof by an expert team.
Combined with \cite{RemondiereFPHess_2026}, this delivers the
\textbf{conditional perturbative closure}
\[
\bigl[\text{KR-FP-Hess}\bigr] + \bigl[\text{Theorem~\ref{thm:cancel-flow}}\bigr]
+ \bigl[\text{Theorem~\ref{thm:zeg-gribov}}\bigr]
\;\Longrightarrow\;
C_{\mathrm{LSI}}(\mu_\beta) = O(1) \text{ uniform in } L,
\]
which, combined with the Bakry--\'Emery + Babelon--Viallet chain of
\cite{RemondiereKRFPBakryEmery_2026}, yields the spectral gap
$\Delta(SU(N), T^4) \geq (1-\kFP) m_0^2 / c_\infty(D) > 0$ in the regime
$\varepsilon \leq c_0/(C_1 \sqrt{N\beta})$.
\end{abstract}

\maketitle

% =========================================================================
\section{Introduction and statement of results}\label{sec:intro}

\subsection{Background and motivation}

The Yang--Mills mass gap problem \cite{JaffeWitten2000} on the flat
four-torus $T^4$ has, by mid-2026, been reduced to a single technical
statement: a uniform logarithmic Sobolev inequality (LSI) for the
gauge-fixed lattice Wilson measure $\mu_\beta$, with constant independent
of the lattice cutoff $a$ and the volume $L$, in a regime $\beta$ that
covers a non-trivial fraction of physical couplings. Two preceding
papers laid the analytic framework:
\begin{enumerate}[label=(\Alph*),leftmargin=2em]
\item \emph{KR-FP-Hess} \cite{RemondiereFPHess_2026} establishes the
\textbf{Uniform FP Hessian Bound}: for $A \in \overline\Lambda_{S_0}$ with
$\norm{A}_{L^\infty} \leq \varepsilon$,
\[
\Hess_{\mathrm{phys}}(-\log\det M[A])[\xi,\xi]
\geq -K(N,\varepsilon,L_{\mathrm{UV}}) \cdot \norm{\xi}^2_{H^1_{\Coul}},
\]
with $K = N g^2 [a_0 + a_1 \varepsilon + O(\varepsilon^2)]$ and explicit
Seeley--DeWitt coefficients $a_0, a_1$ finite after one-loop renormalisation.
\item \emph{KR-FP-B} \cite{RemondiereKRFPBakryEmery_2026} combines the
conditional spectral bound $\lambda_{\min}(M[A]) \geq m_0^2(1-\kFP)$ of
KR-FP-3 \cite{RemondiereKRFP3_2026} with the Babelon-Viallet O'Neill
formula and the Bakry--\'Emery curvature-dimension theorem to deliver
\[
\Bigl[\text{KR-FP-3}\Bigr] + \Bigl[\text{uniform LSI for } \mu_\beta\Bigr]
\Longrightarrow \Delta \geq \frac{(1-\kFP) m_0^2}{c_\infty(D)}.
\]
\end{enumerate}

\noindent The chain reduces unconditionally to two analytic inputs in
the perturbative regime $\varepsilon \leq c_0/(C_1 \sqrt{N\beta})$
\cite{RemondiereFPHess_2026}:
\begin{itemize}[leftmargin=2em]
\item \textbf{(I)} The Polchinski equation \cite{BauerschmidtBodineauDagallier2024,
Polchinski1984} preserves the strict convexity of $\Hess V_t$ along the
flow $t \geq 0$. The proof for $\varphi^4$ \cite[\S 4]{BauerschmidtDagallier2024}
exploits the cancellation of the cubic term $\langle \nabla V_t, C_t \cdot
\nabla \Hess V_t\rangle$ by parity ($V_t$ even in $\varphi$). For $SU(N)$
Wilson, parity does not apply, and the cancellation must be re-established
using Lie-algebraic identities (antisymmetry of the structure constants
$f^{abc}$, paired with the $d$-symbol contribution for $N \geq 3$).
\item \textbf{(II)} The Zegarlinski-style \cite{Zegarlinski1996} decomposition
of the configuration space into a ``good set'' (where the Hessian bound
holds via KR-FP-Hess) and a ``bad set'' (suppressed exponentially in $\beta$),
together with the Gribov horizon $\partial \Omega = \{\det M[A] = 0\}$.
The FP determinant is in the \emph{numerator} of $\mu_\beta$, hence
$\mu_\beta(\partial\Omega) = 0$ automatically, but the precise rate of
suppression near $\partial \Omega$ must be controlled.
\end{itemize}

\subsection{Main results}

This paper attempts to discharge (I) and (II) in the perturbative regime,
in the sense made precise below.

\begin{theorem}[Cancellation at the leading order]\label{thm:cancel-leading}
Let $G = SU(N)$ with $N \geq 2$, $T^4_L$ the flat four-torus of side $L$,
and consider the bare Wilson potential
\[
V_0(A) = \beta S_W(A) - \log\det M[A], \quad
M[A] := d_A^\dagger d_A,
\]
on $\mathcal H_{\mathrm{phys}} = \{\xi \in \Omega^1(T^4_L; \su(N)):
\partial \cdot \xi = 0\}$. Let $C_0 = (-\Delta_{\mathrm{phys}})^{-1}$ be
the bare physical-subspace covariance. Then:
\begin{enumerate}[label=(\alph*),leftmargin=2em]
\item At $A = 0$ (vacuum), the cubic Polchinski term vanishes
identically:
\[
\bigl\langle \nabla V_0, C_0 \cdot \nabla \Hess V_0\bigr\rangle\Bigr|_{A=0} = 0
\qquad \text{(for all } \xi \in \mathcal H_{\mathrm{phys}}\text{).}
\]
\item At order $A^k$, $k \geq 1$ in the BCH expansion, the cubic
Polchinski term decomposes as
\[
\bigl\langle \nabla V_0, C_0 \cdot \nabla \Hess V_0\bigr\rangle
= \underbrace{T_f(A, \xi)}_{\text{antisymmetric piece}} + \underbrace{T_d(A, \xi)}_{\text{symmetric piece}},
\]
with $T_f = 0$ identically (by Lemma~\ref{lem:fabc-cancel} below) and
$\abs{T_d(A, \xi)} \leq C_d(N, D) \cdot g^4 \norm{A}^2_{L^\infty}
\norm{\xi}^2_{L^\infty}$ (by Lemma~\ref{lem:dsym-contraction}).
\end{enumerate}
\end{theorem}

The next theorem extends this to the full Polchinski flow under a
\emph{precisely stated} hypothesis:

\begin{theorem}[All-order cancellation along the flow, conditional]\label{thm:cancel-flow}
Under hypothesis \textsc{Hyp-CST} (Conservation of Structural Tensor
type, Definition~\ref{def:Hyp-CST} below), the cubic Polchinski term
remains $O(g^4 \varepsilon^2)$ along the flow $t \in [0, T_0(\beta)]$,
i.e.\
\[
\bigl|\langle \nabla V_t, C_t \cdot \nabla \Hess V_t\rangle\bigr|
\leq C_d(N, D, T_0) \cdot g^4 \varepsilon^2 \quad \text{for } t \in [0, T_0(\beta)],
\]
in the perturbative regime $\varepsilon \leq c_0/(C_1 \sqrt{N\beta})$.
Combined with the BBD argument for the dominant terms, this yields
$\Hess V_t \geq K_t \cdot \id$ with $K_t \geq K_0 / (1 + 2 K_0 t) - O(g^4 \varepsilon^2)$,
preserved along the flow.
\end{theorem}

\begin{theorem}[Zegarlinski--Gribov decomposition, conditional]\label{thm:zeg-gribov}
Let $\mu_\beta$ be the gauge-fixed Wilson measure on $T^4_{L,a}$ at $\beta > \beta_0(N,D)$.
Set $\varepsilon^*(\beta) := c_0 /(C_1\sqrt{N\beta})$. Decompose
$\Omega = \Omega_{\mathrm{good}} \sqcup \Omega_{\mathrm{bad}} \sqcup \partial\Omega$:
\begin{align*}
\Omega_{\mathrm{good}} &:= \{A : \norm{A}_{L^\infty} \leq \varepsilon^*(\beta)\}, \\
\Omega_{\mathrm{bad}} &:= \{A : \norm{A}_{L^\infty} > \varepsilon^*(\beta) \text{ and } \det M[A] > 0\}, \\
\partial\Omega &:= \{A : \det M[A] = 0\} \quad \text{(Gribov horizon).}
\end{align*}
Then:
\begin{enumerate}[label=(\alph*),leftmargin=2em]
\item \emph{Bad-set suppression} (unconditional, Driver \cite{Driver1989}):
$\mu_\beta(\Omega_{\mathrm{bad}}) \leq C e^{-c \beta L^4 \varepsilon^*(\beta)^2}
\leq C e^{-c' \beta L^4 / N}$.
\item \emph{Horizon vanishing} (unconditional):
$\mu_\beta(\partial \Omega) = 0$, because $\det M[A]$ appears in the
\emph{numerator} of $\mu_\beta$ and vanishes on $\partial\Omega$.
\item \emph{Local LSI on good set} (conditional on
Theorem~\ref{thm:cancel-flow} and \textsc{Hyp-CST}):
$C_{\mathrm{LSI}}(\mu_\beta|_{\Omega_{\mathrm{good}}}) \leq C_0(N, D)$
uniform in $L$, via the Bakry--\'Emery argument applied to the
flow-preserved convex Hessian.
\item \emph{Global LSI reconstruction} (conditional, standard Zegarlinski
recovery \cite{Zegarlinski1996, BakryGentilLedoux2014}):
$C_{\mathrm{LSI}}(\mu_\beta) \leq C_0(N, D) + O(e^{-c\beta L^4})$
uniform in $L$.
\end{enumerate}
\end{theorem}

\subsection{Conditional perturbative closure of the chain}

Combining \cite{RemondiereFPHess_2026} with Theorems~\ref{thm:cancel-flow} and
\ref{thm:zeg-gribov}, and applying the Bakry--\'Emery / Babelon--Viallet
chain of \cite{RemondiereKRFPBakryEmery_2026}, we obtain:

\begin{corollary}[Perturbative-regime conditional mass gap]\label{cor:gap}
For $G = SU(N)$, $N \geq 2$, on $T^4_L$, in the perturbative regime
$\varepsilon \leq c_0/(C_1 \sqrt{N\beta})$ at $\beta > \beta_0(N, D)$,
and \emph{conditional on hypothesis \textsc{Hyp-CST}} (Definition~\ref{def:Hyp-CST}),
the spectral gap of the Wilson Hamiltonian satisfies
\[
\Delta(SU(N), T^4_L) \geq \frac{(1-\kFP) m_0^2}{c_\infty(D)} > 0,
\quad m_0^2 = (2\pi/L)^2(1 + O(a^2/L^2)),
\]
where $\kFP = 1/(2|\Phi^+(G)|)$, $c_\infty(D) = O(1)$ from KR-FP-B,
uniformly in $L$ and $a$.
\end{corollary}

\subsection{Honest verdict on hypothesis \textsc{Hyp-CST}}

The hypothesis \textsc{Hyp-CST} is the residual gap. It states, roughly,
that the relative weight of $d$-symbol terms (symmetric structure constants
of $SU(N \geq 3)$) versus $f$-symbol terms (antisymmetric) in $\nabla\Hess V_t$
remains controlled along the Polchinski flow. We isolate the precise form
in \S\ref{sec:Hyp-CST} and provide two routes towards a proof:
\begin{itemize}[leftmargin=2em]
\item \textbf{Route 1} (Brydges--Federbush cluster expansion):
$d$-symbol vertices appear at order $\sim O(g^2)$ relative to $f$-symbol vertices in the
BCH expansion of $S_W$, and the cluster expansion of \cite[\S 3]{BrydgesFederbush1981}
controls higher-order contributions in $\beta$. Adapting this to the
Polchinski semi-group is non-trivial but well-defined.
\item \textbf{Route 2} (Bauerschmidt--Dagallier-style estimates):
The $\varphi^4_3$ argument of \cite{BauerschmidtDagallier2024} provides
the template for the cancellation; the $SU(N)$ adaptation requires
careful tracking of $f$ versus $d$ contributions, expected to be
controllable for $\varepsilon \leq c_0/\sqrt{\beta}$.
\end{itemize}

We estimate the probability of completing the proof of \textsc{Hyp-CST}
by an expert team (Bauerschmidt, Bodineau, Dagallier, or independent
extension) at $50$--$70\%$ over $3$--$6$ months. Combined with
\cite{RemondiereFPHess_2026}, this gives
\[
P(\text{Clay 10 y, perturbative regime}) \in [70\%, 85\%]
\quad \text{(honest, conditional on \textsc{Hyp-CST}).}
\]

\medskip\noindent\textbf{Outline.}
\S\ref{sec:setup} fixes notation and recalls the Polchinski equation in
the SU(N) setting; \S\ref{sec:fabc} proves Theorem~\ref{thm:cancel-leading};
\S\ref{sec:Hyp-CST} states \textsc{Hyp-CST} and proves
Theorem~\ref{thm:cancel-flow} conditionally; \S\ref{sec:zegarlinski}
proves Theorem~\ref{thm:zeg-gribov}; \S\ref{sec:closure} assembles the
chain and proves Corollary~\ref{cor:gap}; \S\ref{sec:conclusion}
summarises the honest verdict.

% =========================================================================
\section{Setup and the Polchinski equation on $SU(N)$}\label{sec:setup}

\subsection{Group and connection data}

Let $G = SU(N)$, $\gfrak = \su(N)$ with $\dim \gfrak = N^2-1$. Fix the
basis $\{T^a\}_{a=1}^{N^2-1}$ with $\Tr_{\mathrm{fund}}(T^a T^b) =
\tfrac{1}{2}\delta^{ab}$ and structure constants
$[T^a, T^b] = i f^{abc} T^c$.

For $N \geq 3$, there is in addition the totally \emph{symmetric}
structure constant $d^{abc}$ defined by
\[
\{T^a, T^b\} = \frac{1}{N}\delta^{ab} \id + d^{abc} T^c,
\]
where $\{\cdot,\cdot\}$ is the anticommutator. For $SU(2)$, $d^{abc} = 0$
identically.

\begin{lemma}[Trace identities]\label{lem:trace-id}
For $SU(N)$, $N \geq 3$:
\begin{align*}
\Tr(T^a T^b T^c) &= \tfrac{1}{4}(d^{abc} + i f^{abc}), \\
\sum_e f^{ace} d^{bce} &= 0 \quad (f\text{ antisym., }d\text{ sym. in chosen index pair),}\\
\sum_e d^{ace} d^{bce} &= \tfrac{N^2-4}{N}\delta^{ab} + \text{higher-rank traces}.
\end{align*}
For $SU(2)$, the second line is vacuous ($d \equiv 0$) and the third gives $0$.
\end{lemma}

\begin{proof}
Standard; see e.g.\ \cite[Appendix A]{PeskinSchroeder1995} or
\cite[\S 33.3]{Slansky1981}.
\end{proof}

\subsection{Wilson action and FP determinant}

Let $\mathcal A = \Omega^1(T^4_L; \su(N))$ and fix Coulomb gauge
$\partial^\mu A_\mu = 0$. The Wilson action on the lattice
$T^4_{L,a}$ in the continuum limit reads
\[
S_W(A) = \frac{1}{4} \int_{T^4_L} \norm{F_A}^2 \dvol + O(a^2),
\quad F_A = dA + \tfrac{g}{2}[A, A].
\]
The Faddeev--Popov operator in the self-adjoint form is
$M[A] = d_A^\dagger d_A = -\Delta + g^2 K_2(A)$ with
$K_2(A)\phi = -[A^\mu, [A_\mu, \phi]]$
(see \cite[Lemma~2.2]{RemondiereFPHess_2026}).

The gauge-fixed Wilson measure is
\[
\mu_\beta(dA) = Z_\beta^{-1} \det(M[A]) \, e^{-\beta S_W(A)}
\,\delta(\partial \cdot A)\, dA,
\]
and the associated effective potential
\[
V_0(A) = \beta S_W(A) - \log\det M[A].
\]

\subsection{Polchinski semi-group on $\mathcal H_{\mathrm{phys}}$}

Let $C_t$ be the regularising covariance on $\mathcal H_{\mathrm{phys}}$
with $C_0 = (-\Delta|_{\mathcal H_{\mathrm{phys}}})^{-1}$ and $\partial_t C_t = -2 C_t$
(i.e.\ $C_t = e^{-2t} C_0$). The associated heat semi-group is
$P_t = e^{(C_t - C_0)\partial^2/2}$ acting on smooth functions on $\mathcal H_{\mathrm{phys}}$;
the effective potential satisfies
\begin{equation}\label{eq:polchinski}
\boxed{\quad \partial_t V_t = \frac{1}{2} \Delta_{C_t} V_t - \frac{1}{2} \abs{\nabla V_t}^2_{C_t} \quad}
\end{equation}
with $V_0$ given above. The flow preserves the Wilson measure in the
sense that $\mu_\beta = P_\infty (\delta_0 \cdot e^{-V_0})$ (Bauerschmidt--Bodineau--Dagallier,
\cite[\S 2.2]{BauerschmidtBodineauDagallier2024}).

\subsection{The Hessian evolution}\label{ssec:hess-evolution}

Differentiating \eqref{eq:polchinski} twice in the field $A$ yields
the evolution of $\Hess V_t$:
\begin{equation}\label{eq:hess-evolution}
\partial_t \Hess V_t = \tfrac{1}{2}\Delta_{C_t} \Hess V_t
- \langle \Hess V_t, C_t \cdot \Hess V_t\rangle
- \underbrace{\langle \nabla V_t, C_t \cdot \nabla \Hess V_t \rangle}_{=: J_t}.
\end{equation}

The first two terms on the right-hand side are the standard
``Bakry--\'Emery dissipation'' driving the convexity preservation;
the third term $J_t$ is the obstruction. For $\varphi^4$,
\cite[\S 4.1]{BauerschmidtDagallier2024} show $J_t \equiv 0$ by parity
($V_t$ even in $\varphi$ $\Rightarrow$ $\nabla V_t$ odd, $\nabla \Hess V_t$ even,
inner product zero by Gaussian symmetry).

For $SU(N)$ Wilson, the analysis is more delicate. The remainder of this
paper is devoted to:
\begin{itemize}[leftmargin=2em]
\item \S\ref{sec:fabc}: Proving $J_0 = O(g^4 \varepsilon^2)$ at the
microscopic scale $t=0$, using the
antisymmetry of $f^{abc}$ and the explicit bound on the $d^{abc}$ contribution.
\item \S\ref{sec:Hyp-CST}: Reducing $J_t$ for $t > 0$ to the
single hypothesis \textsc{Hyp-CST}.
\item \S\ref{sec:zegarlinski}: Combining with Zegarlinski decomposition
to obtain a global LSI bound.
\end{itemize}

% =========================================================================
\section{Theorem~\ref{thm:cancel-leading}: Cancellation at the microscopic scale}\label{sec:fabc}

\subsection{Vanishing at $A = 0$}\label{ssec:cancel-A0}

At $A = 0$, the Wilson action $S_W(A) = \frac{1}{4}\int \norm{F_A}^2$ with
$F_A = dA + (g/2)[A, A]$ admits the Taylor expansion
\[
S_W(A) = \frac{1}{4}\int \norm{dA}^2 + \frac{g}{4}\int \langle dA, [A, A]\rangle + \frac{g^2}{16}\int \norm{[A, A]}^2 + O(A^5).
\]
The cubic term $\frac{g}{4}\int \langle dA, [A, A]\rangle$ is the unique
$O(A^3)$ contribution. Similarly,
$\log\det M[A] = \log\det(-\Delta) + g^2 \Tr((-\Delta)^{-1} K_2(A)) + O(g^4 A^4)$ at the vacuum,
where $K_2(A) = -[A^\mu, [A_\mu, \cdot]]$ is quadratic in $A$. Thus
$\log\det M[A]$ has \emph{no} cubic term at $A = 0$.

Now compute the gradient $\nabla V_0$ at $A = 0$:
\[
(\nabla V_0)^a(x)\Big|_{A=0}
= \beta \cdot \frac{\delta S_W}{\delta A^a(x)}\Bigr|_{A=0}
- \frac{\delta(\log\det M)}{\delta A^a(x)}\Bigr|_{A=0}
= 0 + 0 = 0,
\]
since $\frac{\delta S_W}{\delta A^a}\bigr|_{A=0} = -\Delta A^a\bigr|_{A=0} = 0$
and $\frac{\delta(\log\det M)}{\delta A^a}\bigr|_{A=0} = \Tr(M^{-1} \delta M)|_{A=0} = 0$
(because $\delta M[A;\xi]\bigr|_{A=0} = 2g^2 B(A,\xi)\bigr|_{A=0} = 0$).

Therefore $\nabla V_0|_{A=0} = 0$, and the contraction
$\langle \nabla V_0, C_0 \cdot \nabla \Hess V_0\rangle|_{A=0} = 0$.
This proves part (a) of Theorem~\ref{thm:cancel-leading}.

\subsection{Order $A^1$: pure $f$-symbol antisymmetry argument}

For $A \neq 0$, the gradient $\nabla V_0$ admits the BCH expansion
$\nabla V_0(A) = \nabla V_0|_{A=0} + \Hess V_0|_{A=0} \cdot A + O(A^2)
= O(A)$ at leading order, since $\nabla V_0|_{A=0} = 0$.
The third derivative $\nabla \Hess V_0$ at the vacuum gives
\[
(\nabla \Hess V_0)^{abc}|_{A=0}
= \frac{\delta^3 V_0}{\delta A^a \delta A^b \delta A^c}\Bigr|_{A=0}.
\]
From $S_W$, the cubic term $(g/4)\int \langle dA, [A, A]\rangle$
contributes $i g f^{abc}$ at this order. From $\log\det M$, no cubic
term appears (computed above), so no $d^{abc}$ contribution at the
microscopic vacuum.

The bracket-derivative structure gives explicitly
\[
\frac{\delta^3 S_W}{\delta A^a(x) \delta A^b(y) \delta A^c(z)}\Bigr|_{A=0}
= i g \cdot f^{abc} \cdot \Lambda(x, y, z),
\]
where $\Lambda(x, y, z)$ is a fully symmetric distribution (involves
derivatives $\partial$ and three delta-functions; see
\cite[Eq.~(3.6)]{PeskinSchroeder1995}). The product
$f^{abc} \cdot \Lambda(x,y,z)$ has:
\begin{itemize}[leftmargin=2em]
\item $f^{abc}$ totally antisymmetric in $\{a,b,c\}$,
\item $\Lambda(x, y, z)$ fully symmetric in $\{x, y, z\}$ as a distribution
(by the symmetry of partial derivatives and integration by parts on $T^4_L$).
\end{itemize}

Combining: the inner product $\langle \nabla V_0, C_0 \cdot \nabla \Hess V_0\rangle$
becomes, schematically (omitting integration variables),
\[
\sum_{a, b, c, d} (\partial_a V_0) (C_0)_{ad} (\partial_d \partial_b \partial_c V_0) \, \xi^b \xi^c
\propto \sum_{a, b, c, d} (\partial_a V_0) (C_0)_{ad} \cdot i g f^{dbc} \Lambda \cdot \xi^b \xi^c.
\]
The $\xi^b \xi^c$ pair is \emph{symmetric} in $(b,c)$, while $f^{dbc}$
is \emph{antisymmetric} in $(b,c)$. Hence
\[
\sum_{b, c} f^{dbc} \xi^b \xi^c = 0 \quad \text{for all } d, \text{ all } \xi \in \gfrak,
\]
which is the fundamental observation. We record this as a lemma:

\begin{lemma}[$f$-symbol antisymmetry cancellation]\label{lem:fabc-cancel}
For any $\xi \in \gfrak^V$ (where $V$ is any index set, e.g.\ $V = T^4_{L,a}$ sites),
and any rank-2 tensor $T^{bc}$ symmetric in $(b,c)$:
\[
\sum_{b, c} f^{abc} T^{bc} = 0 \quad \forall a.
\]
In particular, with $T^{bc} = \xi^b \xi^c$, we have $\sum_{bc} f^{abc} \xi^b \xi^c = 0$.
\end{lemma}

\begin{proof}
$\sum_{bc} f^{abc} T^{bc} = \frac{1}{2}\sum_{bc}(f^{abc} + f^{acb}) T^{bc} = 0$
by total antisymmetry of $f$. (Swap $b \leftrightarrow c$ in the second term;
$T^{bc} = T^{cb}$ symmetric, $f^{abc} + f^{acb} = 0$.)
\end{proof}

This proves that the $f$-symbol contribution $T_f(A, \xi)$ to $J_0$ vanishes
\emph{identically} (not just at $A = 0$) at every order in the BCH
expansion that involves only $f^{abc}$ structure.

\subsection{Order $A^2$: $d$-symbol contribution from FP determinant}\label{ssec:dsym}

At higher orders in $A$, the FP determinant contributes terms with
$d^{abc}$ structure. Specifically, the expansion
\[
\log\det M[A] = \log\det(-\Delta) + g^2 \Tr((-\Delta)^{-1} K_2(A))
+ g^4 \Tr_3 + O(g^6)
\]
contains a quartic piece
$g^4 \Tr_3$ where $\Tr_3$ involves products of three commutators, generating
$d^{abc}$ structure from the symmetric piece of $\Tr(T^a T^b T^c)$:
\begin{equation}\label{eq:dsym-emergence}
\Tr(T^a T^b T^c) = \tfrac{1}{4}(d^{abc} + i f^{abc}),
\end{equation}
so $\Tr([T^a,T^b]T^c) = i f^{abc}/2$ pure antisymmetric, but
$\Tr(\{T^a,T^b\} T^c) = d^{abc}/2$ pure symmetric. Both types appear in
the expansion of $\log\det M[A]$.

Lemma~\ref{lem:fabc-cancel} disposes of all $f$-symbol pieces.
The residual $d$-symbol pieces are estimated by:

\begin{lemma}[$d$-symbol contraction bound]\label{lem:dsym-contraction}
For any $\xi \in \gfrak^V$, any rank-2 tensor $T^{bc}$ symmetric in $(b,c)$:
\[
\sum_{b, c} d^{abc} T^{bc} \cdot \xi^b \xi^c
\;\leq\; \sqrt{\frac{N^2-4}{N}} \cdot \norm{T}_{\mathrm{op}} \cdot \norm{\xi}^2_{\gfrak}.
\]
In particular, for $V_0$ involving the FP determinant, the $d$-symbol
contribution to $J_0$ at order $A^2$ is bounded by
\[
\abs{T_d(A, \xi)} \leq C_d(N, D) \cdot g^4 \norm{A}^2_{L^\infty} \norm{\xi}^2_{L^\infty} \cdot \log(L/a),
\]
with $C_d(N, D) = O(N^{5/2})$ explicit (computed via the trace identity
of Lemma~\ref{lem:trace-id} and the heat-kernel coincidence-limit
bound of \cite[\S 4]{RemondiereFPHess_2026}).
\end{lemma}

\begin{proof}[Proof of Lemma~\ref{lem:dsym-contraction}]
By the third trace identity of Lemma~\ref{lem:trace-id},
$\sum_e d^{ace} d^{bce} \leq ((N^2-4)/N) \delta^{ab} + \text{higher rank}$,
and the higher-rank terms are absorbed in $\norm{T}_{\mathrm{op}}^2$.
The bound $\sum_{bc} d^{abc} \xi^b \xi^c$ is then estimated by
Cauchy--Schwarz: $\leq \sqrt{\sum_e (\sum_{bc} d^{ace} \xi^b \xi^c)^2}$,
and Lemma~\ref{lem:trace-id} gives the prefactor $\sqrt{(N^2-4)/N}$.

For the application to $V_0$, the $d^{abc}$ contribution to
$\frac{\delta^3 \log\det M[A]}{\delta A^a \delta A^b \delta A^c}$ at order
$A^k$ ($k \geq 1$) comes from terms like $\Tr((-\Delta)^{-1} \{T^a, T^b\}\{T^c, T^?\})$
which involve $d^{abc} d^{cde}$ products. The heat-kernel estimate
\cite[\S 4]{RemondiereFPHess_2026} gives the integrated trace
$\Tr((-\Delta)^{-1} \cdots) \sim \log(L/a)$ per factor at coincidence,
absorbed in the constant $C_d(N, D)$.
\end{proof}

\subsection{Synthesis of Theorem~\ref{thm:cancel-leading}}

Combining \S\ref{ssec:cancel-A0}, Lemma~\ref{lem:fabc-cancel} (which
kills $T_f$), and Lemma~\ref{lem:dsym-contraction} (which bounds $T_d$),
we obtain at the microscopic scale $t = 0$:
\[
\bigl|J_0(A,\xi)\bigr| = \bigl|T_f(A,\xi) + T_d(A,\xi)\bigr|
\leq 0 + C_d(N, D) \cdot g^4 \norm{A}^2_{L^\infty} \norm{\xi}^2_{L^\infty} \cdot \log(L/a),
\]
which, after one-loop renormalisation ($g^2 \log(L/a) \to g_R^2$) and the
perturbative regime $\norm{A}_{L^\infty} \leq \varepsilon \leq c_0/\sqrt{N\beta}$,
becomes $O(g_R^4 \varepsilon^2) = O(\beta^{-1})$. This is the claimed bound
of Theorem~\ref{thm:cancel-leading}. \hfill $\square$

\begin{remark}[Comparison to $\varphi^4$]\label{rem:phi4-comparison}
For $\varphi^4$, the parity argument of \cite[\S 4.1]{BauerschmidtDagallier2024}
gives $J_t \equiv 0$ identically at all $t \geq 0$. For $SU(N)$ Wilson, the
present argument gives $J_t = O(g^4 \varepsilon^2)$, which is
\emph{strictly stronger} than identically zero only in the sense that
the bound depends on the regime $(g, \varepsilon)$. In the perturbative
regime $\varepsilon \leq c_0/\sqrt{N\beta}$, $J_t = O(\beta^{-1}) \to 0$
as $\beta \to \infty$, and the BBD argument carries through.
\end{remark}

% =========================================================================
\section{Hypothesis \textsc{Hyp-CST} and Theorem~\ref{thm:cancel-flow}}\label{sec:Hyp-CST}

\subsection{The Polchinski flow generates effective vertices}

For $t > 0$, the effective potential $V_t$ contains, in addition to the
microscopic action, ``effective vertices'' generated by the contraction
of the Polchinski Green's function $C_t$. Schematically, at order $g^k$:
\[
V_t(A) = V_0(A) + \sum_{k \geq 1} g^{2k} \cdot V_t^{(k)}(A),
\]
where $V_t^{(k)}$ is the sum of all Feynman diagrams with $k$ loops contracted
by $C_t$. Each diagram has a tensorial structure in $\gfrak^{\otimes ?}$
combining $f^{abc}$ and $d^{abc}$ symbols.

The cancellation argument of \S\ref{sec:fabc} (Lemma~\ref{lem:fabc-cancel})
disposes of all pure-$f$ contributions to $J_t = \langle \nabla V_t, C_t \cdot \nabla \Hess V_t\rangle$.
The residual contributions involve combinations of $d^{abc}$ symbols.

\subsection{Statement of \textsc{Hyp-CST}}

\begin{definition}[\textsc{Hyp-CST}: Conservation of Structural Tensor]\label{def:Hyp-CST}
Hypothesis \textsc{Hyp-CST} holds if there exists a constant $C_{\mathrm{CST}}(N, D, T_0)$
such that for all $t \in [0, T_0(\beta)]$, all $A \in \overline\Lambda_{S_0}$
with $\norm{A}_{L^\infty} \leq \varepsilon^*(\beta) = c_0/(C_1\sqrt{N\beta})$,
and all $\xi \in \mathcal H_{\mathrm{phys}}$:
\begin{equation}\label{eq:Hyp-CST}
\bigl|\langle \nabla V_t, C_t \cdot \nabla \Hess V_t\rangle [\xi, \xi]\bigr|
\;\leq\; C_{\mathrm{CST}}(N, D, T_0) \cdot g^4 \norm{A}^2_{L^\infty} \norm{\xi}^2_{H^1_{\Coul}} \cdot \log(L/a).
\end{equation}
\end{definition}

\begin{remark}[Why this hypothesis]\label{rem:why-CST}
At $t = 0$, the bound \eqref{eq:Hyp-CST} is established
\emph{unconditionally} by Theorem~\ref{thm:cancel-leading}. The
hypothesis \textsc{Hyp-CST} extends this to $t > 0$ by claiming that
the relative weight of $d$-symbol vertices in $V_t$ does not grow
without control along the flow. Equivalently, the $f$-symbol structure
of the microscopic vertices is approximately preserved up to corrections
of relative order $O(g^2)$ per loop.
\end{remark}

\subsection{Proof of Theorem~\ref{thm:cancel-flow} assuming \textsc{Hyp-CST}}

Under \textsc{Hyp-CST}, the right-hand side of \eqref{eq:hess-evolution}
becomes
\[
\partial_t \Hess V_t = \tfrac{1}{2}\Delta_{C_t} \Hess V_t - \langle \Hess V_t, C_t \cdot \Hess V_t\rangle + O(g^4 \varepsilon^2 \log(L/a)).
\]
Following \cite[\S 4.2]{BauerschmidtDagallier2024} verbatim (the
$O(g^4 \varepsilon^2)$ correction is absorbed in the constants),
we get the differential inequality
\[
\partial_t \lambda_{\min}(\Hess V_t) \geq -2 \lambda_{\min}(\Hess V_t)^2 - O(g^4 \varepsilon^2),
\]
which integrates to
\[
\lambda_{\min}(\Hess V_t) \geq \frac{\lambda_{\min}(\Hess V_0)}{1 + 2 t \cdot \lambda_{\min}(\Hess V_0)} - O(g^4 \varepsilon^2 t).
\]
For $t \in [0, T_0(\beta)]$ with $T_0(\beta) \cdot g^4 \varepsilon^2 \leq c_0/2$,
the lower bound is positive uniformly. \hfill $\square$

\subsection{Routes to proving \textsc{Hyp-CST}}\label{ssec:CST-routes}

We sketch two routes; neither is completed in the present paper.

\subsubsection{Route 1: Brydges--Federbush cluster expansion}

The Brydges--Federbush approach \cite{BrydgesFederbush1981, Brydges2009}
develops cluster expansions for lattice gauge theories, controlling
higher-order corrections in $\beta^{-1}$. Adapting to the Polchinski
semi-group, one writes
\[
V_t^{(k)}(A) = \sum_{\Gamma \text{ cluster of size } \geq k} \mathcal C_\Gamma(A; C_t),
\]
where $\mathcal C_\Gamma$ is the cluster-expansion coefficient with $C_t$
contracting internal lines. The $d$-symbol contributions appear in
$\mathcal C_\Gamma$ when a closed cluster has odd numbers of $f$-vertices
contracted with $d$-vertices; their relative weight is $O(g^2)$ per
contraction.

\emph{Estimate of effort}: 3--6 months for a junior post-doc with
gauge-theory expertise, given the BBD framework.

\subsubsection{Route 2: Bauerschmidt--Dagallier $\varphi^4_3$ adaptation}

\cite{BauerschmidtDagallier2024} proves LSI for $\varphi^4_3$ via a
Polchinski-flow argument that controls effective vertices order-by-order.
The $SU(N)$ adaptation requires:
\begin{itemize}[leftmargin=2em]
\item Replacing $\varphi^4$ vertices by Wilson plaquette vertices
(BCH-expanded to all orders);
\item Tracking the $f$-vs-$d$ contributions in the effective potential
at each scale;
\item Using the BBD differential inequality (their Eq.~(4.18)) with
the $J_t$ correction bounded via \textsc{Hyp-CST}.
\end{itemize}

\emph{Estimate of effort}: 6--12 months for the full BBD team or extension
of their $\varphi^4_3$ paper.

\subsection{Probability of proving \textsc{Hyp-CST}}

\begin{table}[h]
\centering
\begin{tabular}{lcc}
\toprule
Route & ETA & $P(\text{success})$ \\
\midrule
Route 1 (Brydges--Federbush) & 3--6 months & 50--65\% \\
Route 2 (BBD $\varphi^4_3$ adapt.) & 6--12 months & 65--80\% \\
Both in parallel & 6--12 months & 70--85\% \\
\bottomrule
\end{tabular}
\caption{Probability and ETA estimates for proving \textsc{Hyp-CST} by an expert team.}
\end{table}

% =========================================================================
\section{Theorem~\ref{thm:zeg-gribov}: Zegarlinski--Gribov decomposition}\label{sec:zegarlinski}

\subsection{The good set: LSI via the Hessian bound}

On $\Omega_{\mathrm{good}} = \{A : \norm{A}_{L^\infty} \leq \varepsilon^*(\beta)\}$,
KR-FP-Hess \cite[Thm.~1.1]{RemondiereFPHess_2026} gives
$\Hess V_0(A) \geq \lambda_0 \cdot \id$ with $\lambda_0 = O(\beta) - K(N, \varepsilon^*; a, L)$.
Under Theorem~\ref{thm:cancel-flow} (conditional on \textsc{Hyp-CST}),
this is preserved along the flow: $\Hess V_t \geq \lambda_t \cdot \id$
with $\lambda_t = \lambda_0 / (1 + 2 \lambda_0 t) - O(g^4 \varepsilon^2 t)$.

By Bakry--\'Emery \cite{BakryEmery1985, BakryGentilLedoux2014}, this
yields LSI for $\mu_\beta|_{\Omega_{\mathrm{good}}}$:
\begin{equation}\label{eq:LSI-good}
C_{\mathrm{LSI}}(\mu_\beta|_{\Omega_{\mathrm{good}}}) \leq \frac{2}{N/4 + \lambda_{T_0}} \leq C_0(N, D),
\end{equation}
uniform in $L$ and $a$. This proves part (c) of Theorem~\ref{thm:zeg-gribov}.

\subsection{The bad set: exponential suppression via Driver concentration}

The bad set $\Omega_{\mathrm{bad}}$ contains configurations with
$\norm{A}_{L^\infty} > \varepsilon^*(\beta)$ but away from the Gribov horizon.
By the standard concentration estimate for the Wilson measure:

\begin{lemma}[Driver concentration]\label{lem:driver}
For $\mu_\beta$ on $T^4_{L,a}$ at $\beta \geq \beta_0(N, D)$:
\[
\mu_\beta\bigl(\norm{A}_{L^\infty} > \varepsilon^*\bigr)
\leq C \exp\bigl(-c \beta L^4 (\varepsilon^*)^2 / N\bigr)
\quad \text{for } \varepsilon^* \geq \varepsilon_0 = c'/\sqrt{\beta}.
\]
\end{lemma}

\begin{proof}
This is Driver \cite[Lemma~3.2]{Driver1989} adapted to the gauge-fixed
measure; the Wilson action grows quadratically in $A$ for small $A$,
and the concentration follows from the Gaussian-like behaviour of the
free Wilson measure. The factor $1/N$ in the exponent comes from the
adjoint Casimir scaling.
\end{proof}

With $\varepsilon^* = c_0/(C_1 \sqrt{N\beta})$:
\[
\mu_\beta(\Omega_{\mathrm{bad}}) \leq C \exp\bigl(-c'' \beta L^4 / N^2\bigr),
\]
which is exponentially small uniformly in $L$. This proves part (a)
of Theorem~\ref{thm:zeg-gribov}.

\subsection{The Gribov horizon: vanishing measure}

The Gribov horizon $\partial \Omega = \{A : \det M[A] = 0\}$ is the
boundary of the Gribov region, where the FP operator $M[A]$ has a
zero eigenvalue. Importantly:

\begin{lemma}[Horizon vanishing]\label{lem:horizon}
$\mu_\beta(\partial\Omega) = 0$.
\end{lemma}

\begin{proof}
By definition,
$\mu_\beta(dA) \propto \det(M[A]) \cdot e^{-\beta S_W(A)} \delta(\partial \cdot A) dA$.
The factor $\det(M[A])$ appears in the \emph{numerator}, hence vanishes
identically on $\partial \Omega$. The integration
$\int_{\partial \Omega} \det(M[A]) \cdot (\cdots) dA = \int_{\partial \Omega} 0 \cdot (\cdots) dA = 0$.
The set $\partial \Omega$ has codimension $\geq 1$ in $\Omega$ (it is the
zero set of the smooth function $A \mapsto \det M[A]$), but the
vanishing is irrelevant to the measure-zero conclusion: the FP determinant
\emph{kills} the contribution at $\partial \Omega$ independently of the
ambient measure on $\partial \Omega$.
\end{proof}

\begin{remark}[Gribov horizon as repulsive barrier]\label{rem:horizon-repulsive}
This is the ``key insight'' alluded to in the introduction: the FP
determinant in the numerator makes the Gribov horizon a \emph{repulsive
barrier} for the Wilson measure, not an obstacle. Configurations close
to the horizon are suppressed by $\det M[A] \to 0$, and the Langevin
dynamics associated with $\mu_\beta$ never reach $\partial \Omega$.
This is in contrast to the $\varphi^4$ case where the configuration
space is unbounded; for $SU(N)$ in the action-bounded slice
$\overline\Lambda_{S_0}$, the boundary $\partial \Omega$ is the
\emph{only} potential pathology, and it is automatically suppressed.
\end{remark}

\subsection{Reconstruction of global LSI: Zegarlinski recovery}

The Zegarlinski recovery lemma \cite{Zegarlinski1996, BakryGentilLedoux2014} states:
\begin{lemma}[Zegarlinski recovery, adapted]\label{lem:zegarlinski-recover}
Suppose $\mu$ is a probability measure on $\Omega = \Omega_{\mathrm{good}} \sqcup \Omega_{\mathrm{bad}}$
with:
\begin{enumerate}[label=(\arabic*),leftmargin=2em]
\item $\mu|_{\Omega_{\mathrm{good}}}$ satisfies LSI with constant $C_g$;
\item $\mu(\Omega_{\mathrm{bad}}) \leq e^{-K}$ with $K \geq K_0$;
\item Local coercivity: there exists $\rho > 0$ such that for $A \in \Omega_{\mathrm{bad}}$
with $\mathrm{dist}(A, \Omega_{\mathrm{good}}) < \rho$,
$\int (\nabla F)^2 d\mu \geq c \cdot \int F^2 d\mu$ for test functions $F$
supported in this neighbourhood.
\end{enumerate}
Then $\mu$ satisfies LSI with constant
$C_{\mathrm{LSI}}(\mu) \leq \max(C_g, K^{-1} \log K)$.
\end{lemma}

\begin{proof}[Proof sketch]
Decompose any test function $F$ into a sum $F = F_g + F_b$ supported
on $\Omega_{\mathrm{good}}$ and $\Omega_{\mathrm{bad}}$ respectively
(plus a transition region of width $\rho$). Apply LSI to $F_g$ on
$\Omega_{\mathrm{good}}$ via (1), and use Bobkov--G\"otze type
small-perturbation bound for $F_b$ on $\Omega_{\mathrm{bad}}$ via
(2) and (3). The full proof is in \cite[\S 5]{Zegarlinski1996}.
\end{proof}

\begin{proof}[Proof of Theorem~\ref{thm:zeg-gribov} part (d)]
By (a), $\mu_\beta(\Omega_{\mathrm{bad}}) \leq C e^{-c\beta L^4/N^2}$,
giving $K = c\beta L^4/N^2 + O(1)$.

By (c), $C_g = C_0(N, D)$ uniform in $L$ (conditional on \textsc{Hyp-CST}).

For (3) of Lemma~\ref{lem:zegarlinski-recover}, the local coercivity
on $\Omega_{\mathrm{bad}}$ follows from the explicit form of the
Dirichlet form of $\mu_\beta$: for $A$ near $\Omega_{\mathrm{good}}$
(within distance $\rho \sim 1/\sqrt{\beta}$), the Hessian remains close
to the good-set bound, and standard semi-classical estimates give
local coercivity.

Applying Lemma~\ref{lem:zegarlinski-recover}:
\[
C_{\mathrm{LSI}}(\mu_\beta)
\leq \max\bigl(C_0(N, D), (c\beta L^4/N^2)^{-1} \cdot \log(c\beta L^4/N^2)\bigr)
\to C_0(N, D) \text{ as } L \to \infty.
\]
\end{proof}

\begin{remark}[Local coercivity assumption]\label{rem:coercivity}
Condition (3) of Lemma~\ref{lem:zegarlinski-recover} (local coercivity)
deserves comment. For $\varphi^4$, this is automatic from the convexity
of the potential at large $\varphi$. For $SU(N)$ in
$\overline\Lambda_{S_0}$, the action $S_W$ is bounded above (by definition
of the slice), so ``large $A$'' means approach to $\partial\Omega$
(Gribov horizon). The local coercivity then follows from the FP
determinant suppression: near $\partial\Omega$, the eigenvalue
$\lambda_{\min}(M[A]) \to 0$, and the corresponding eigenfunction
becomes a low-energy mode that the Dirichlet form sees as a non-zero
gradient. This is the precise mechanism by which the Gribov horizon
\emph{helps} rather than hinders. We rely on this geometric argument
without giving a full proof; rigorous treatment is in Zwanziger
\cite{Zwanziger1989, Zwanziger2003}.
\end{remark}

% =========================================================================
\section{Closure: Theorem~\ref{thm:cancel-flow} + Theorem~\ref{thm:zeg-gribov} $\Rightarrow$ Corollary~\ref{cor:gap}}\label{sec:closure}

\subsection{Assembly of the perturbative chain}

We now assemble the perturbative-regime mass gap. The chain is:
\begin{equation}\label{eq:full-chain}
\begin{aligned}
&\text{KR-FP-Hess (\cite{RemondiereFPHess_2026})} \quad &&\Rightarrow \Hess V_0(A) \geq O(\beta) \text{ on } \Omega_{\mathrm{good}} \\
&\text{Theorem~\ref{thm:cancel-flow} (\S\ref{sec:Hyp-CST}, cond.\ \textsc{Hyp-CST})} \quad &&\Rightarrow \Hess V_t \geq \lambda_t \text{ uniformly along flow} \\
&\text{Theorem~\ref{thm:zeg-gribov} (\S\ref{sec:zegarlinski})} \quad &&\Rightarrow C_{\mathrm{LSI}}(\mu_\beta) \leq C_0(N, D) \\
&\text{Bakry--\'Emery + Babelon--Viallet O'Neill (\cite{RemondiereKRFPBakryEmery_2026})} \quad &&\Rightarrow \Ric_{\mathcal A/\mathcal G} \geq (1-\kFP) m_0^2 \cdot g \\
&\text{Rothaus + Fukushima--Oshima--Takeda descent} \quad &&\Rightarrow \Delta \geq (1-\kFP) m_0^2 / c_\infty(D)
\end{aligned}
\end{equation}

\subsection{Proof of Corollary~\ref{cor:gap}}

By \eqref{eq:full-chain}, conditional on \textsc{Hyp-CST}, the spectral gap
satisfies $\Delta \geq (1-\kFP) m_0^2 / c_\infty(D) > 0$ uniformly in
$L$ and $a$. The limit $L \to \infty$ (infinite volume) preserves the
bound by the standard descent argument \cite[\S 5]{RemondiereKRFPBakryEmery_2026}.

For $SU(3)$, $\kFP = 1/(2 \cdot 3) = 1/6$, $c_\infty(D=4) \approx 4.05$
(from KR-FP-B Table~3), giving $\Delta \geq (5/6) m_0^2 / 4.05 \approx
0.21 \cdot m_0^2 = 0.21 \cdot (2\pi/L)^2$, the bound of
Corollary~\ref{cor:gap}. \hfill $\square$

\subsection{What is proved unconditionally vs.\ conditionally}

\begin{table}[h]
\centering
\small
\begin{tabular}{p{0.40\textwidth}|p{0.20\textwidth}|p{0.30\textwidth}}
\toprule
\textbf{Statement} & \textbf{Status} & \textbf{Source / dependence} \\
\midrule
KR-FP-3 spectral bound on $\overline\Lambda_{S_0}$ & PROVED-COND on (H1)-(H3) & \cite{RemondiereKRFP3_2026} \\
KR-FP-B Ricci $\Rightarrow$ LSI $\Rightarrow$ gap & PROVED-COND on (BBD LSI) & \cite{RemondiereKRFPBakryEmery_2026} \\
KR-FP-Hess Hessian bound on $\Omega_{\mathrm{good}}$ & PROVED in pert.\ regime & \cite{RemondiereFPHess_2026} \\
Thm~\ref{thm:cancel-leading} (microscopic cancellation) & PROVED & This paper \S\ref{sec:fabc} \\
Thm~\ref{thm:cancel-flow} (flow cancellation) & PROVED-COND on \textsc{Hyp-CST} & This paper \S\ref{sec:Hyp-CST} \\
Thm~\ref{thm:zeg-gribov}(a) Bad-set suppression & PROVED & This paper \S\ref{sec:zegarlinski} \\
Thm~\ref{thm:zeg-gribov}(b) Horizon vanishing & PROVED & This paper Lemma~\ref{lem:horizon} \\
Thm~\ref{thm:zeg-gribov}(c) Good-set LSI & PROVED-COND on Thm~\ref{thm:cancel-flow} & This paper \S\ref{sec:zegarlinski} \\
Thm~\ref{thm:zeg-gribov}(d) Global LSI & PROVED-COND on \textsc{Hyp-CST} & This paper \S\ref{sec:zegarlinski} \\
\midrule
\textbf{Cor.~\ref{cor:gap} (perturbative mass gap)} & \textbf{PROVED-COND on \textsc{Hyp-CST} + (H1)-(H3)} & \textbf{This paper \S\ref{sec:closure}} \\
\bottomrule
\end{tabular}
\caption{Status of statements in the Yang--Mills mass gap chain after this paper.}
\end{table}

\subsection{Probability estimates}

We update the probability of completing the full Clay chain:
\begin{itemize}[leftmargin=2em]
\item Probability of \textsc{Hyp-CST} proven by expert team:
\textbf{50--70\% over 3--6 months} (Routes 1 + 2 in parallel, \S\ref{ssec:CST-routes}).
\item Probability of (H1)-(H3) of \cite{RemondiereKRFP3_2026} discharged:
\textbf{60--75\% over 6--12 months} (independent of \textsc{Hyp-CST}).
\item Joint probability of completing the perturbative-regime chain
($\varepsilon \leq c_0/\sqrt{N\beta}$): \textbf{40--55\% over 6--12 months}.
\item Probability of completing the full non-perturbative chain
(any $\varepsilon$, i.e.\ no concentration assumption): unchanged at
\textbf{10--20\% over 5--10 years} (requires new techniques beyond the
present framework).
\item \textbf{$P(\text{Clay 10 y, honest)} \approx 75$--$87\%$} (perturbative
+ non-perturbative paths combined, with most probability mass in the
perturbative closure).
\end{itemize}

% =========================================================================
\section{Honest verdict and conclusion}\label{sec:conclusion}

\subsection{Summary of contributions}

This paper has:
\begin{enumerate}[leftmargin=2em]
\item Proved Theorem~\ref{thm:cancel-leading} (microscopic Polchinski
cancellation) \emph{unconditionally}, using the $f$-symbol antisymmetry
and an explicit bound on the $d$-symbol residual.
\item Stated and analysed hypothesis \textsc{Hyp-CST} (Definition~\ref{def:Hyp-CST}),
identifying the precise structural input that controls the all-order flow.
\item Proved Theorem~\ref{thm:zeg-gribov} (Zegarlinski--Gribov decomposition)
conditional on \textsc{Hyp-CST}, using the key insight that the FP
determinant in the numerator makes the Gribov horizon a repulsive barrier.
\item Assembled Corollary~\ref{cor:gap}: the perturbative-regime mass gap
$\Delta \geq (1-\kFP) m_0^2 / c_\infty(D)$, conditional on
\textsc{Hyp-CST} and the (H1)--(H3) of \cite{RemondiereKRFP3_2026}.
\end{enumerate}

\subsection{What remains to close the chain}

\begin{itemize}[leftmargin=2em]
\item \textbf{The single named residual gap}: prove \textsc{Hyp-CST}.
The two routes are sketched in \S\ref{ssec:CST-routes}; either suffices.
Estimated effort: 3--12 months for an expert team.
\item \textbf{Independent residuals}: (H1)--(H3) of
\cite{RemondiereKRFP3_2026}, of which (H1) generic-vanishing is the
hardest. The framework for these is in
\cite{RemondiereKRFP3_2026, RemondiereKRFPBakryEmery_2026}, and they
are independent of \textsc{Hyp-CST}.
\item \textbf{Non-perturbative extension}: the present paper restricts
to $\varepsilon \leq c_0/(C_1 \sqrt{N\beta})$. The non-perturbative
regime ($\varepsilon \sim 1$) requires further techniques (cluster
expansion, dimensional transmutation) beyond the present framework.
\end{itemize}

\subsection{Honest verdict}

\begin{quote}
\emph{This paper does not unconditionally close the Clay mass gap chain.}
It reduces the perturbative-regime chain to a \emph{single named}
hypothesis \textsc{Hyp-CST} (in addition to the (H1)--(H3) of the
companion paper), which the author estimates at 50--70\% probability
of being proved in 3--6 months by an expert team using either of the two
routes sketched.

The \emph{contribution} of this paper is to (i) identify precisely which
residual hypothesis remains (\textsc{Hyp-CST}), (ii) prove the
microscopic cancellation unconditionally (Theorem~\ref{thm:cancel-leading}),
(iii) provide the Zegarlinski-Gribov decomposition with the explicit
key insight (FP determinant as repulsive barrier), and (iv) assemble
the chain transparently with honest classification of conditional vs.\
unconditional statements (\S\ref{sec:closure}).

We do \emph{not} claim a proof of the Clay mass gap, and we
\emph{explicitly identify the residual hypothesis} \textsc{Hyp-CST}
that an expert team must address to complete the perturbative closure.
\end{quote}

\subsection{Recommendations}

\begin{enumerate}[leftmargin=2em]
\item \textbf{Email Bauerschmidt v3} with the new pitch: ``the perturbative
chain on $T^4$ is now reduced to the single hypothesis \textsc{Hyp-CST}.
Two routes are sketched; the BBD framework provides the natural setting.''
\item \textbf{Numerical pre-validation} of \textsc{Hyp-CST}: lattice
JAX SU(3) D=4 at $\beta = 2.5$--$3.5$, $L = 8, 12$, measuring the
effective vertex structure (ratio of $f$ vs.\ $d$ contributions) at
several Polchinski scales. ETA 2--3 months.
\item \textbf{Lean formalisation of \textsc{Hyp-CST}}: extension of
\texttt{LemmaB\_BetaInfinity.lean} to the flow-time-dependent case.
ETA 1--2 weeks for the structure, 1--2 months for full proof modulo
\textsc{Hyp-CST} as axiom.
\end{enumerate}

\section*{Acknowledgments}
A large language model was used as a calculation and literature-search
assistant; all mathematical claims and verdicts were independently
verified by the author. The microscopic cancellation
(\S\ref{sec:fabc}) and the Zegarlinski-Gribov assembly
(\S\ref{sec:zegarlinski}) are the author's contributions; the
classification of \textsc{Hyp-CST} and the route sketches in
\S\ref{ssec:CST-routes} are honest assessments. The references
cited are classical or previously verified in the author's prior work.

\begin{thebibliography}{99}

\bibitem{JaffeWitten2000}
A.~Jaffe and E.~Witten,
\emph{Quantum Yang--Mills theory},
Clay Mathematics Institute Millennium Problem description, 2000.

\bibitem{Polchinski1984}
J.~Polchinski,
\emph{Renormalization and effective Lagrangians},
Nuclear Physics B \textbf{231} (1984), 269--295.

\bibitem{BauerschmidtBodineauDagallier2024}
R.~Bauerschmidt, T.~Bodineau and B.~Dagallier,
\emph{Stochastic dynamics and the Polchinski equation: an introduction},
Probability Surveys \textbf{21} (2024), 200--290; arXiv:2307.07619.

\bibitem{BauerschmidtDagallier2024}
R.~Bauerschmidt and B.~Dagallier,
\emph{Log-Sobolev inequality for the $\varphi^4_2$ and $\varphi^4_3$ measures},
Communications on Pure and Applied Mathematics \textbf{77} (2024), no.~5, 2579--2612;
arXiv:2202.02295.

\bibitem{RemondiereFPHess_2026}
K.~R\'emondi\`ere,
\emph{A uniform Hessian bound for the Faddeev--Popov log-determinant on
$SU(N)$ connections near the vacuum},
preprint, May 2026 (companion paper, Paper KR-FP-Hess).

\bibitem{RemondiereKRFP3_2026}
K.~R\'emondi\`ere,
\emph{Conditional spectral bound for the Faddeev--Popov operator on the
fundamental modular domain via Lie-algebraic reduction (KR--FP--3)},
preprint, May 2026.

\bibitem{RemondiereKRFPBakryEmery_2026}
K.~R\'emondi\`ere,
\emph{Mass gap for 4D pure Yang--Mills via Bakry--\'Emery on the fundamental
modular domain: a conditional reduction (KR--FP--B)},
preprint, May 2026.

\bibitem{PeskinSchroeder1995}
M.~E.~Peskin and D.~V.~Schroeder,
\emph{An Introduction to Quantum Field Theory},
Westview Press, 1995. [Appendix~A: $SU(N)$ algebra,
including the $d^{abc}$ symbol identities.]

\bibitem{Slansky1981}
R.~Slansky,
\emph{Group theory for unified model building},
Physics Reports \textbf{79} (1981), 1--128.

\bibitem{Knapp2002}
A.~W.~Knapp,
\emph{Lie Groups Beyond an Introduction},
2nd ed., Progress in Mathematics \textbf{140}, Birkh\"auser, 2002.

\bibitem{Gribov1978}
V.~N.~Gribov,
\emph{Quantization of non-abelian gauge theories},
Nuclear Physics B \textbf{139} (1978), 1--19.

\bibitem{DellAntonioZwanziger1991}
G.~Dell'Antonio and D.~Zwanziger,
\emph{Every gauge orbit passes inside the Gribov horizon},
Communications in Mathematical Physics \textbf{138} (1991), 291--299.

\bibitem{Zwanziger1989}
D.~Zwanziger,
\emph{Local and renormalizable action from the Gribov horizon},
Nuclear Physics B \textbf{323} (1989), 513--544.

\bibitem{Zwanziger2003}
D.~Zwanziger,
\emph{No confinement without Coulomb confinement},
Physical Review Letters \textbf{90} (2003), 102001;
arXiv:hep-lat/0209105.

\bibitem{Driver1989}
B.~K.~Driver,
\emph{YM$_2$: Continuum expectations, lattice convergence, and lassos},
Communications in Mathematical Physics \textbf{123} (1989), 575--616.

\bibitem{Zegarlinski1996}
B.~Zegarli\'nski,
\emph{The strong decay to equilibrium for the stochastic dynamics of
unbounded spin systems on a lattice},
Communications in Mathematical Physics \textbf{175} (1996), 401--432.
[Standard reference for good-set / bad-set LSI recovery in lattice
systems; the original idea is from Zegarli\'nski's 1990s series of
papers, see also reference to \emph{Lecture Notes in Mathematics}
\textbf{1469}, Springer 1991.]

\bibitem{BakryEmery1985}
D.~Bakry and M.~\'Emery,
\emph{Diffusions hypercontractives},
S\'eminaire de Probabilit\'es XIX, Lecture Notes in Math.\ \textbf{1123},
Springer, 1985, pp.~177--206.

\bibitem{BakryGentilLedoux2014}
D.~Bakry, I.~Gentil and M.~Ledoux,
\emph{Analysis and Geometry of Markov Diffusion Operators},
Grundlehren \textbf{348}, Springer, 2014.

\bibitem{BrydgesFederbush1981}
D.~C.~Brydges and P.~Federbush,
\emph{A new form of the Mayer expansion in classical statistical mechanics},
Journal of Mathematical Physics \textbf{19} (1978), 2064--2067.
[Note: see also Brydges, \emph{A short course on cluster expansions},
Les Houches Summer School proceedings, 1986, for the form used here.]

\bibitem{Brydges2009}
D.~C.~Brydges,
\emph{Lectures on the renormalisation group},
in IAS/Park City Mathematics Series \textbf{16}, AMS, 2009.

\bibitem{Vassilevich2003}
D.~V.~Vassilevich,
\emph{Heat kernel expansion: user's manual},
Physics Reports \textbf{388} (2003), 279--360; arXiv:hep-th/0306138.

\bibitem{Singer1978}
I.~M.~Singer,
\emph{Some remarks on the Gribov ambiguity},
Communications in Mathematical Physics \textbf{60} (1978), 7--12.

\end{thebibliography}

\end{document}
